% from main.tex
%%%%%%%%%%  scan idx 102  |  folio 89  %%%%%%%%%%
En
fait, les exemples b) et c) sont des cas particuliers de l'exemple suivant :

\textbf{d)} Soient \((S,\mathcal O_S)\) un topos annelé (cf. SGA 4 IV 2.10)\footnote{\(\mathcal O_S\) est ici un anneau de \(S\), à élément unité, mais non nécessairement commutatif.}, et \(\mathcal C\) la \(S\)-catégorie scindée définie par \(U\longmapsto\) (catégorie des \(\mathcal O_U\)-Modules à gauche)~; \(\mathcal C\) est une \(S\)-catégorie abélienne plate. Soit \(\mathcal C_0\) la sous-\(S\)-catégorie additive strictement pleine de \(\mathcal C\) définie par \(U\longmapsto\) (catégorie des \(\mathcal O_U\)-Modules libres de type fini). Un \(\mathcal O_S\)-Module localement de \(\mathcal C_0\)-type fini sera dit simplement de \emph{type fini}. La catégorie \(\mathcal C_0\) est localement quasi-stable par noyau d'épimorphisme et localement relevable dans \(\mathcal C\), comme il résulte du

\textbf{Lemme 1.3.1.} \emph{Soit \(P\) un \(\mathcal O_S\)-Module, Conditions équivalentes} :
\begin{enumerate}[label=(\roman*)]
\item \emph{\(P\) est de type fini et le foncteur \(\Hom(P,\cdot)\) est exact};
\item \emph{\(P\) est de type fini et tout épimorphisme \(M\to P\) admet localement une section};
\item \emph{\(P\) est localement facteur direct d'un Module libre de type fini}.
\end{enumerate}
La preuve est immédiate et laissée en exercice au lecteur.

\subsection*{1.4. Un lemme de «~prolongement~» (cf. EGA \(0_{\mathrm{III}}\) 11.9.1)}

\textbf{Proposition 1.4.1.} Soient \(S\) un site, \(\mathcal C\) une \(S\)-catégorie abélienne plate (1.1), \(\mathcal C_0\) une sous-\(S\)-catégorie additive (1.1) strictement pleine de \(\mathcal C\); on suppose \(\mathcal C_0\) localement quasi-relevable dans \(\mathcal C\) (1.2). Soient \(S\in\Ob S\), \(u:E\longrightarrow F\) un morphisme (de degré \(0\)) de complexes de \(\mathcal C_S\) (différentielles de degré \(+1\)), \(G\) le cône de \(u\) (cf. [V] chap.\ I \S~1 2.1), \(n\in\mathbb Z\); on suppose que \(\Ho^n(G)\) est de \(\mathcal C_0\)-
